% HD-G16. 21 November 2026. GPT-6 Astra Pro, inherited effort.
\section{Unequal scalar twists: an explicit isogeny comparison}

\paragraph{HD-G16.}
The whole-character obstruction of HD-G15 also holds for
the actual object
\[
 \mathcal E_{r,t}
   =\Phi\bigl((I_0\otimes\mathcal O_X(rH))
                 \boxtimes(I_0\otimes\mathcal O_X(tH))\bigr),
 \qquad r,t\in\mathbb Z,
                                                               \tag{G16.1}
\]
with the same CM fourfold and Fourier equivalence as(G15.5).
The proof gives an actual isogeny relating the cohomological
calculation to the equal-twist calculation.

Put
\[
 \ell=\frac{r+t}{2},\qquad \delta=\frac{r-t}{2},\qquad
 P'=P-2\delta A,\qquad D'=D-\delta P+\delta^2A.
                                                               \tag{G16.2}
\]
These are stated rational cohomology classes. Write
$v_{r,t}=e^{-P/2}\operatorname{ch}(\mathcal E_{r,t})$.
The general version of the first integral in(G15.8) is
\[
 \int_{c,d}e^{r h_{x+y}+t h_y-P_{\mathrm{in}}}
  =(r+t)\exp\left(\frac{rt}{r+t}ab
                    +\frac{r}{r+t}P_{\mathrm{out}}
                    -\frac{1}{r+t}uv\right)
                                                               \tag{G16.3}
\]
when $r+t\ne0$. After multiplication by
$e^{-P_{\mathrm{out}}/2}$, its polynomial form, valid for
every $r,t$, is
\[
 (r+t)+rt\,ab+\delta P_{\mathrm{out}}-uv
                         -\frac{r+t}{4}abuv.               \tag{G16.4}
\]
Substituting(G16.2) turns this exactly into the equal-twist
polynomial(G15.9) with parameter $\ell$ and second-factor
form $D'$.
The other two exponential terms become
\[
        e^{tA+P/2}=e^{\ell A+P'/2},\qquad
        e^{rA-P/2}=e^{\ell A-P'/2},
\]
and the point term is unchanged. Taking the four products
therefore gives precisely(G15.6), with
$(\ell,A,D,P)$ replaced by $(\ell,A,D',P')$.
For $\ell=0$, the same polynomial calculation gives
(G15.7) in the primed forms.
Its rank is
\[
                         (r+t)^4-2\ne0.                  \tag{G16.5}
\]
The nonvanishing holds because $r+t$ is an integer.

Here is the actual morphism implementing the primed
coordinate comparison. Identify $\widehat X$ with $X$
using the specified principal polarization and define
\[
 S:X\times X\longrightarrow X\times X,\qquad
            S(x,\beta)=(2x,\;2\beta+(r-t)x).
                                                               \tag{G16.6}
\]
Its integral matrix has nonzero determinant, so it is
an isogeny. Its degree is $2^{16}$, as follows by taking
the real lattice determinant of its two diagonal
multiplication-by-two blocks. The addition identity
for the product principal polarization gives, directly,
\[
                       S^*A=4A,\qquad
                       S^*P'=4P,\qquad
                       S^*D'=4D.                         \tag{G16.7}
\]
For example
$S^*P=4P+8\delta A$ and
$S^*D=4D+4\delta P+4\delta^2A$, which verify all three
identities without an implicit coordinate normalization.

Consequently the pullback of the centered second class
$C(v_{r,t})$ is sixteen times the centered second class
in the equal-twist calculation at the rational parameter
$\ell$. The proof of HD-G15 only used nonzero rank for
that calculation; (G16.5) supplies it here.
Thus $S^*C(v_{r,t})$ is invariant under the actual
involution $\operatorname{id}\times[-1]$ and has
nonzero component in
$H^{2,0}(X)\otimes H^{0,2}(X)$.
The explicit coefficient $1/8$ there becomes $2$.

For a whole character of the form(G15.1), pullback by
the actual isogeny $S$ would give
$S^*C=\gamma(S^*\theta)^2$ with $S^*\theta$ ample.
The positive-square uniqueness and mixed-component
argument of HD-G15 therefore give the same contradiction.
No claim that $S$ already preserves a chosen integral
$K$-action is required: the argument uses this actual
pullback of the character and polarization, and the
Weil term continues to have codimension four.

This proves the obstruction for all unequal integer
scalar twists in(G16.1). It does not assert it for
arbitrary line-bundle classes or arbitrary other
Fourier constructions. \hfill$\square$
